\section{Validation: Cartel Adjacency Under Coarsened Observability}
\label{sec:cade}

The validation answers two questions in sequence. First: does
the screening statistic concentrate on populations with a higher
prior probability of operating inside cartel-affected
environments than random matching at comparable participation
volume? Second: does the concentration appear in evidence that
predates estimation, so that the result is not an artifact of
the within-sample CADE record? We answer both in the
affirmative, against the structurally appropriate validation
object given the data layer.

That object is cartel \emph{adjacency}, not membership. The
award-layer envelope (§\ref{sec:institutional}) does not contain
the information required for firm-level membership
identification, and the screening framing of
§\ref{sec:emp_estimand} does not require it: the participation
primitive the framework identifies is precisely a loser-side
exposure measure, not a winner-side label. The cobidder
population---always-loser firms that participated alongside CADE
direct defendants---is therefore the empirical content the
screening statistic is designed to recover. We report against
direct-defendant labels separately to bound the construct's
scope and to document the design's empirical signature.

\subsection{Validation populations}
\label{sec:cade_objects}

The validation reads against five nested populations, each playing
a distinct role in the empirical exercise.
Table~\ref{tab:cade_populations} catalogs them, with sample sizes
and the function of each in the validation. The leftmost column
moves from the broadest universe (all BEC participants) inward
through the always-loser stratum to the flagged frequent-loser
subset; the two rightmost rows split the CADE adjudication
record into direct defendants---typically frequent winners and
therefore outside the construct's target---and cobidders inside
the always-loser stratum, which are the cartel-adjacency labels
the screening statistic is designed to recover.

\begin{table}[htbp]
\centering
\caption{Validation populations and object discipline}
\label{tab:cade_populations}
\footnotesize
\begin{adjustbox}{max width=\textwidth,center}
\begin{tabular}{p{4.0cm}p{2.0cm}p{6.5cm}}
\toprule
Population & Size & Definition and role \\
\midrule
All BEC participants & $41{,}444$ & Universe of firms bidding in BEC, $2009$--$2019$. \\
Always-losers & $\valAlwaysLosers$ & Firms with zero wins across the sample; stratum within which the screening statistic operates. \\
Frequent losers & $\valFL$ & Always-losers above the median $+\,1.5\!\times$IQR participation threshold; the flagged population. \\
CADE direct defendants in BEC & $\valDirectCADE$ & Firms named in CADE rulings; typically frequent winners; outside the screening statistic's target population. \\
Cobidders in the always-loser stratum & $\valCobidders$ & Always-losers that participated alongside CADE direct defendants; the cartel-adjacency label the screening statistic is designed to recover. \\
\bottomrule
\end{tabular}
\end{adjustbox}

\smallskip
{\footnotesize\textit{Notes:} The screening statistic operates
within the always-loser stratum, not across the broader BEC
universe. Participation alongside an adjudicated cartel member
is consistent with cover bidding, capacity-transfer motives,
market intelligence, or bystander participation in
cartel-affected markets. Strong AUC against this label
corroborates that the screening statistic concentrates on
cartel-relevant environments; it does not adjudicate
membership.\par}
\end{table}

\subsection{Within-stratum profile of cobidders}
\label{sec:cade_bridge}

The validation positive class is cobidders inside the always-loser
stratum, not direct CADE defendants. The reading would be
unconvincing if cobidders were just always-loser firms with a
proximity label glued on. Table~\ref{tab:theory_bridge} compares
cobidders against three reference populations along four
operational predictions of the cover-bidder type that the framework
of Online Appendix~A formalizes.

% Subprompt 4: theory-validation bridge (script 60)
\begin{table}[htbp]
\centering
\caption{Cobidders vs.\ Reference Populations: Operational Predictions of the Cover-Bidder Type}
\label{tab:theory_bridge}
\begin{threeparttable}
\footnotesize
\begin{adjustbox}{max width=\textwidth}
\begin{tabular}{lcccccc}
\toprule
 & Cobidders & FL non- & Always- & Other & Cohen's $d$ & Wilcoxon \\
 & ($N=191$) & cobidders & loser, non-FL & winners & vs.\ FL & $p$ \\
\midrule
Mean portfolio item-group HHI & $0.380$ & $0.288$ & $0.719$ & $0.294$ & $+0.39$ & $<0.001$ \\
Mean share of loss-bids facing a direct CADE defendant & $0.015$ & $0.002$ & $0.002$ & $0.003$ & $+0.46$ & $<0.001$ \\
Mean share of loss-bids in winner-pairs of count $\geq 5$ & $0.216$ & $0.334$ & $0.026$ & $0.412$ & $-0.38$ & $<0.001$ \\
Mean total participations & $136.471$ & $76.652$ & $12.119$ & $1603.320$ & $+0.67$ & $<0.001$ \\
\bottomrule
\end{tabular}
\end{adjustbox}
\begin{tablenotes}
\footnotesize
\item \textit{Notes:} Each row reports a firm-level mean for four populations: cobidders ($\valCobidders$ always-losers participating alongside CADE direct defendants), FL non-cobidders ($\valFL{-}\!\valCobidders$ frequent losers without recorded direct-defendant adjacency), always-losers below the FL threshold, and other winners (firms with at least one win, excluding direct CADE defendants). The two rightmost columns report Cohen's $d$ and a Wilcoxon rank-sum test of cobidders against FL non-cobidders, the relevant within-stratum comparison. The four operational measures map directly to the cover-bidder type in Online Appendix~A: total participation (deployment intensity), share of own loss-bids facing a direct CADE defendant (deployment proximity), share of loss-bids in winner-pairs of count $\geq 5$ (repeat co-bidding density), and portfolio item-group HHI (specialization in the cartel's market).
\item \textit{Source:} \texttt{scripts/60\_theory\_validation\_bridge.R}.
\end{tablenotes}
\end{threeparttable}
\end{table}


The within-stratum comparison---cobidders versus FL non-cobidders,
both above the FL participation threshold---moves in the direction
the framework predicts on $4$ of the $5$ measured dimensions where
the data layer can speak. Cobidders deploy at higher intensity:
$\valBridgeTendCob$ mean total participations against
$\valBridgeTendFLnc$ for FL non-cobidders, Cohen's $d=\valBridgeTendD$,
$p<10^{-15}$; the gap is sharper still on unique winners faced
($\valBridgeUniqWinCob$ vs $\valBridgeUniqWinFLnc$,
$d=\valBridgeUniqWinD$, $p<10^{-22}$). They are deployed proximally
to the cartel: $\valBridgeDirectCob$ of their loss-bids face a direct
CADE defendant against $\valBridgeDirectFLnc$ for FL non-cobidders,
$d=\valBridgeDirectD$, $p<10^{-70}$. They specialize in narrower
market slices: portfolio item-group HHI of $\valBridgeHHICob$
against $\valBridgeHHIFLnc$ ($d=\valBridgeHHID$,
$p<10^{-13}$), touching $\valBridgeNGroupsCob$ distinct item-groups
on average against $\valBridgeNGroupsFLnc$
($d=\valBridgeNGroupsD$, $p<10^{-6}$). The one dimension that
moves the wrong way is the share of own loss-bids in repeat-pair
clusters of count $\geq 5$
($\valBridgeRepeatShareCob$ vs $\valBridgeRepeatShareFLnc$,
$d=\valBridgeRepeatShareD$); we note it without leaning on it.

The contrast with non-FL always-losers is sharper still. Firms
below the FL threshold participate in approximately $12$ tenders
each, face $2$ unique winners on average, touch only $2$ item-groups,
and almost never face a direct CADE defendant: a thin-participation,
narrow-portfolio profile that does not look like the cover-bidder
type. The cobidder population is empirically distinct from this
reference group as well as from FL non-cobidders.

The bridge is partial. The framework also predicts firm-level bid
aggressiveness and bid dispersion patterns; the award-record data
layer does not preserve per-firm bid amounts and these predictions
cannot be tested directly here. This is a portability point: in
jurisdictions where bid microdata is preserved at the firm level,
the bridge could be extended in those two further dimensions
(\S\ref{sec:limitations}, \S\ref{sec:forensic}). Within what the
data can measure, cobidders match the modeled type along four of
five operational predictions, including the two with the largest
standardized differences. We read the discrimination evidence in
the rest of this section against an empirically grounded validation
object, not against a label of convenience.

\subsection{Conservative benchmark}
\label{sec:cade_contemporaneous}

The primary benchmark restricts the ground truth to
$\valConservativeCases$ CADE cases adjudicated before
$\valConservativeCutoff$ (coleta de lixo, aquecedores solares,
trens-metr\^os, tecnologia-informa\c{c}\~ao $2019$-$10$):
$\valConservativeFD$ firm-defendants, all active in BEC, with
$\valConservativeCobidders$ always-losers co-bidding with at
least one of them, of which $\valConservativeFL$ are frequent
losers.

The relevant comparison is whether frequent losers co-bid with
direct defendants at rates above what random matching at
comparable participation volume would produce. We construct the
counterfactual baseline by stratified permutation: for each of
$\valPermPilot$ iterations, we permute the cobidder/non-cobidder
labels within strata defined by participation-volume quartile,
re-compute the share of frequent losers among the relabeled
cobidders, and accumulate the empirical distribution under the
null of random matching conditional on participation volume.
The participation-stratified co-bidding rate observed in the
data is $\valCADEenrich$ against a permutation baseline of
$\valCADEbase$ ($p<0.001$, $\valPermPilot$ iterations), an
excess ratio of $\valMultThreeTwo\!\!\times$. The ROC AUC
against the contemporaneous ground truth is $\valAUCprePost$
($95\%$ DeLong CI $[0.713, 0.783]$).

This is the primary benchmark for the screening interpretation
on three grounds. The construct uses no information from cases
adjudicated post-sample. The labels predate estimation: cartel
existence and membership were established by enforcement before
the screening statistic was built. And the permutation baseline
strips out the mechanical co-bidding rate that participation
volume would produce by itself, isolating the residual
concentration the screening statistic is meant to identify. The
$\valMultThreeTwo\!\!\times$ excess is what the screening
framing predicts under the separating equilibrium of
Online Appendix~A: where the cartel deploys cover bidders, the
loser-side participation footprint concentrates above the
random-matching baseline.

\subsection{Prospective benchmark and within-firm enrichment}
\label{sec:cade_prospective}

A more permissive benchmark uses CADE's full procurement-cartel
portfolio ($12$ cases, $65$ firm-defendants, of which
$\valDirectCADE$ are active in BEC). Among $\valFL$ frequent
losers, $\valCobidders$ ($\valCobidShareFL$) co-participate with
at least one of the $\valDirectCADE$; the
participation-stratified excess ratio is
$\valMultThreeFive\!\!\times$ ($p<0.001$). Firm-level AUC is
$\valAUCFLfirm$ ($95\%$ CI $\valAUCFLfirmCI$) in-sample and
$\valAUCFLfirmTemp$ ($\valAUCFLfirmTempCI$) under temporal
holdout (train $2009$--$2016$, test $2017$--$2019$). The
benchmark treats post-sample adjudications as informative labels
for cartel involvement during the BEC window---a reasonable but
non-trivial assumption that the conservative benchmark in
\S\ref{sec:cade_contemporaneous} avoids. We report it as
complementary, not as the primary anchor.

A within-firm exercise on the same portfolio: $3$ of $7$
always-loser direct defendants ($\valPctFortythree$) cross the
frequent-loser threshold against a population baseline of
$\valPctSixteenRef$. Table~\ref{tab:cade_fl_firms} lists them.

\begin{table}[htbp]
\centering
\caption{Direct CADE defendants flagged by the construct}
\label{tab:cade_fl_firms}
\begin{adjustbox}{max width=\textwidth,center}
\small
\begin{tabular}{p{3.5cm}p{4.0cm}lcc}
\toprule
Firm & Cartel & CADE process & Tenders & Wins \\
\midrule
Sol Tecnologia & Solar water heaters & 08012.001273/2010-24 & 84 & 0 \\
Nova Esperan\c{c}a Locadora & School transportation & 08700.005876/2019-85 & 97 & 0 \\
Jofran Com\'ercio & Trash bags (Op.\ Colludium) & 08700.005789/2015-02 & 65 & 0 \\
\bottomrule
\end{tabular}
\end{adjustbox}
\end{table}

These are the exception within a direct-defendant population
that is otherwise frequent-winner-heavy; the AUC against direct
defendants in the broader BEC firm universe remains
$\valAUCdirectStd$. The asymmetry between cobidder
discrimination ($\valAUCFLfirm$ in-sample,
$\valAUCFLfirmTemp$ under temporal holdout) and direct-defendant
discrimination ($\valAUCdirectStd$) is the design's empirical
signature, not a failure of validation. The screening statistic
recovers loser-side participation footprints; direct defendants
are by construction the winner side of the same arrangements
and lie outside the target population.

\subsection{Robustness to the enforcement record}
\label{sec:cade_robustness}

Dropping all $\valCADEpermPerm$ CADE-involved tender-items and
re-estimating the within-PBU baseline yields
$\hat\beta = \valExclCADECoef$ ($N = \valCADEpermObs{,}954$),
virtually unchanged from $\valPctSixfour$. The broad-sample
screening signal therefore does not depend on the within-sample
CADE adjudications for its content: removing the items where
the cartel record bites leaves the conditional association
intact. Under the screening framing this is the predicted
outcome---the participation footprint the construct flags is a
loser-side equilibrium object, present wherever the deployment
problem operates, and not a feature mechanically produced by the
items already labeled by enforcement. The robustness completes
the conservative benchmark of §\ref{sec:cade_contemporaneous}:
the labels that drive the validation are stripped out of the
estimation sample, and the price-imprint signal survives.

\subsection{What the validation does and does not show}
\label{sec:cade_scope}

The validation corroborates the screening interpretation along
three dimensions, each consistent with the framework of
Online Appendix~A. The participation-stratified excess of
$\valMultThreeTwo\!\!\times$ on the conservative benchmark, the
discrimination AUC of $\valAUCprePost$ on cases adjudicated
before the sample window closes, and the
$\valMultThreeFive\!\!\times$ excess on the full prospective
portfolio (AUC $\valAUCFLfirm$ in-sample, $\valAUCFLfirmTemp$
under temporal holdout) all read in the same direction:
within the always-loser stratum, frequent losers concentrate
near adjudicated cartel members at rates well above
participation-stratified random matching.

Three boundaries delimit what the validation supports. The
labels are indirect by construction, a feature of working at
the award layer rather than the bid layer; the AUC asymmetry
between cobidders and direct defendants
($\valAUCFLfirm$ versus $\valAUCdirectStd$) is the empirical
signature of a screening statistic that recovers loser-side
participation, not winner-side identity; and the screening
signal in $\beta$ persists when CADE-involved items are removed
from estimation, which means the validation labels are not the
mechanism that produces the price-imprint result.

The validation establishes that the screening statistic flags
cartel-adjacent environments. Under coarsened observability it
cannot adjudicate firm-level membership, and the screening
framing does not require it to.
